\scanpage{55}
\deleted{\unsure{E.|E}} Cela prouvera bien 1.3. car pour prouver $S_X({}^{\deleted{\unsure{t|*}}}Z(\unsure{\varpi|\mathfrak{a}|\mathfrak{A}}))=u_Z(S_Y(\varpi))$ \inserted{\unsure{dans|sur} $A_X(k)$}, il suffit de prouver (grâce à la \deleted{bi}dualité grossière i.e. ensembliste \deleted{$A_X$ \unsure{sur}} déduite de $A_X\xrightarrow{\sim}A_X^{**}=B^*$, \unsure{à savoir} $A_X(k)\longrightarrow\mathrm{Pic}^0(B)$) que $\Theta(S_X({}^{\deleted{\unsure{t|*}}}Z(\varpi)))=\Theta(u_Z(S_Y(\varpi)))$ \inserted{dans $\mathrm{Pic}^0(B)$}. \deleted{\unsure{Or par const.}} Cela résultera de 1.4. et \inserted{de} l'équation
\[
\deleted{\Theta(u_Z(b))=c'(b)\quad\inserted{\text{dans }\deleted{\unsure{B}}\mathrm{Pic}^0(B_X)}\quad\text{pour tout } b\in A_Y(k)}
\]
\[
c'(\varpi)=\Theta(u_Z(S_Y(\varpi)))\qquad\text{dans }\mathrm{Pic}^0(B)
\]
qui résulte de la définition de $u_Z$ en termes de $c'$, cf préliminaires.

D'autre part, 1.4. sera une conséquence immédiate du

\uline{Lemme} 1.5. Soit $\varpi$ comme dans 1.4. Alors il existe un diviseur $D$ dans la classe \unsure{d'|de'}équivalence linéaire de $\Theta$, \unsure{telle|tell|tel} que, \deleted{\unsure{l'on ait}} posant $C=\pi_*(Z''\cdot(D\times Y))\in\deleted{\unsure{Z}}\mathrm{Cyc}^1(B\times Y)$, \deleted{\illegible{1}} on ait
\[
D(\varphi_{X*}({}^{\deleted{\unsure{t|*}}t}Z(\varpi)))=\deleted{\unsure{Z^t|{}^t Z}}C(\varpi)\qquad\text{dans }\mathrm{Cyc}^1(B)
\]
(\inserted{et que} les deux membres de cette égalité soient définis !).
% archive sheet number: 54
